\subsection{\texorpdfstring{$L^p$}{Lp} Spaces}

% **Role:** modern integration and almost-everywhere analysis.
%
% Lp spaces, H\"older, Minkowski, completeness, density, duality, Riesz representation.

Lebesgue spaces turn integrability into normed linear geometry. The essential point is that an integral norm cannot distinguish functions that differ only on null sets, so the actual vectors are almost-everywhere equivalence classes. Their bounded linear functionals are themselves represented by integration against functions in the conjugate space.

% Source guide: The Elements of Integration, Chapter 6, Sections 6.1--6.16.
\begin{definition} \label{def:norms-and-seminorms-lp-source} \hlt{Norms and Seminorms}\\
Let $\mathcal{V}$ be a vector space over $\mathbb{F}\in\{\R,\C\}$.
\begin{enumerate}[label=(\roman*)]
\setlength{\itemsep}{0pt}
\item A \hlt{seminorm} is $N:\mathcal{V}\to[0,\infty)$ with $N(\alpha v)=\abs{\alpha}N(v)$ and $N(u+v)\leq N(u)+N(v)$ $\forall u,v\in\mathcal{V}$, $\alpha\in\mathbb{F}$.
\item A \hlt{norm} is a seminorm with $N(v)=0\Longleftrightarrow v=0$. A norm $N$ induces $d_N(u,v)\defeq N(u-v)$.
\end{enumerate}
\end{definition}

\begin{example} \label{ex:coordinate-sequence-function-norms-seminorms} \hlt{Norms and Seminorms Before Quotienting}
\begin{enumerate}[label=(\roman*)]
\setlength{\itemsep}{0pt}
\item On $\mathbb{F}^n$, $\norm{u}_p\defeq(\sum\limits_{j=1}^n\abs{u_j}^p)^{\frac{1}{p}}$ for $1\leq p<\infty$ and $\norm{u}_\infty\defeq\max\limits_{1\leq j\leq n}\abs{u_j}$ are norms.\\
Minkowski gives $\norm{u+v}_p\leq\norm{u}_p+\norm{v}_p$.
\item On sequences, $\ell^p\defeq\{a:\sum\limits_n\abs{a_n}^p<+\infty\}$ with $\norm{a}_p\defeq(\sum\limits_n\abs{a_n}^p)^{\frac{1}{p}}$.\\
$\ell^\infty\defeq\{a:\sup\limits_n\abs{a_n}<+\infty\}$ with $\norm{a}_\infty\defeq\sup\limits_n\abs{a_n}$.
\item On $B(X)$, $\norm{f}_\infty\defeq\sup\limits_{x\in X}\abs{f(x)}$ is a norm. On $C[0,1]$, $N_0(f)\defeq\abs{f(0)}$ and $N_1(f)\defeq\sup\limits_{0\leq x\leq\frac{1}{2}}\abs{f(x)}$ are seminorms. On $C^1[a,b]$, $N_D(f)\defeq\sup\limits_{a\leq x\leq b}\abs{f'(x)}$ is a seminorm with kernel the constant functions.
\end{enumerate}
\hlt{Key Insight}: Seminorms become norms after quotienting out their zero directions.
\end{example}

\begin{remark} \label{rem:continuous-functions-integral-metric-crossref} \hlt{Integral Norm on Continuous Functions}\\
On $C[0,1]$, $N(f)\defeq\int_0^1\abs{f(x)}\,dx$ is a norm: if $f(x_0)\neq0$, continuity gives $N(f)>0$. It is not complete, since continuous ramps can be Cauchy in $N$ and converge in $L^1$ to a discontinuous step function.
\end{remark}

\begin{definition} \label{def:lp-and-linfty-lebesgue-spaces} \hlt{$L^p$ Spaces and Mean Convergence}\\
Let $(X,\mathcal{M},\mu)$ be a measure space.
\begin{enumerate}[label=(\roman*)]
\setlength{\itemsep}{0pt}
\item For $1\leq p<\infty$, the \hlt{$p$-integrable class} is $\mathscr{L}^p(\mu)\defeq\{f:X\to\mathbb{F}:f\text{ measurable},\ \int_X\abs{f}^p\,d\mu<+\infty\}$.
\item The \hlt{$L^p$ space} and \hlt{$L^p$ norm} are $L^p(\mu)\defeq\mathscr{L}^p(\mu)/{=_{\mathrm{a.e.}}}$ and $\norm{f}_p\defeq(\int_X\abs{f}^p\,d\mu)^{\frac{1}{p}}$.
\item \hlt{$L^p$ convergence}, or \hlt{mean convergence of order $p$}, means $f_n\to f\Longleftrightarrow\norm{f_n-f}_p\to0$.\\
\hlt{$L^p$-Cauchy} means $\norm{f_m-f_n}_p\to0$ as $m,n\to\infty$.
\item \hlt{Essentially bounded} means $\abs{f}\leq M$ a.e. for some $M<+\infty$; \hlt{$L^\infty(\mu)$} is the a.e.-quotient of such functions, with \hlt{essential sup norm} $\norm{f}_\infty\defeq\inf\{M\geq0:\abs{f}\leq M\text{ a.e.}\}=\inf\limits_{\mu(N)=0}\sup\limits_{x\in X\setminus N}\abs{f(x)}$.
\end{enumerate}
\end{definition}

\begin{proposition} \label{prop:l1-seminorm-and-ae-quotient} \hlt{$L^1$ as a Quotient of a Seminorm}\\
On $\mathscr{L}^1(\mu)$ set $N_\mu(f)\defeq\int_X\abs{f}\,d\mu$.
\begin{enumerate}[label=(\roman*)]
\setlength{\itemsep}{0pt}
\item $N_\mu$ is a seminorm.
\item $N_\mu(f)=0\Longleftrightarrow f=0$ a.e.
\item $L^1(\mu)\defeq\mathscr{L}^1(\mu)/{=_{\mathrm{a.e.}}}$ is normed by $\norm{[f]}_1\defeq\int_X\abs{f}\,d\mu$; write $f$ for $[f]$.
\end{enumerate}
\begin{proof}
\begin{enumerate}[label=(\roman*)]
\setlength{\itemsep}{0pt}
\item[$(i)$] If $f,g\in\mathscr{L}^1(\mu)$ and $a,b\in\R$, then $\abs{af+bg}\leq\abs a\abs f+\abs b\abs g\in\mathscr{L}^1(\mu)$, so $\mathscr{L}^1(\mu)$ is linear. Also $N_\mu(af)=\abs aN_\mu(f)$ and $N_\mu(f+g)\leq N_\mu(f)+N_\mu(g)$.
\item[$(ii)$] ($\Rightarrow$) If $N_\mu(f)=0$, set $E_n\defeq\{\abs f\geq\frac{1}{n}\}$. Then $\frac{1}{n}\mu(E_n)\leq\int_{E_n}\abs f\,d\mu\leq N_\mu(f)=0$, $\{f\neq0\}=\bigcup\limits_nE_n$ null. 
\item[$(ii)$] ($\Leftarrow$) If $f=0$ a.e., then $\abs f=0$ a.e., hence $N_\mu(f)=0$.
\item[$(iii)$] If $f=g$ a.e., then $\abs f=\abs g$ a.e. and $\int_X\abs f\,d\mu=\int_X\abs g\,d\mu$, so $\norm{[f]}_1$ is well defined. The seminorm axioms descend from (i), and $\norm{[f]}_1=0\Rightarrow f=0$ a.e. by (ii), hence $[f]=[0]$.
\end{enumerate}
\hlt{Structural Commentary}: The seminorm kernel is exactly the a.e.-zero functions.\\
\hlt{Key Insight}: Quotient by null differences before calling the integral a norm.
\end{proof}
\end{proposition}

\begin{definition} \label{def:complex-integrable-and-l2} \hlt{Complex Integrals and $L^2$}
\begin{enumerate}[label=(\roman*)]
\setlength{\itemsep}{0pt}
\item For $f=u+iv$, \hlt{complex measurability} means measurability of $u,v$.
\item On $E\in\mathcal{M}$, \hlt{complex integrability} means $u,v\in\mathscr{L}^1(\mu)$ on $E$; set $\int_Ef\,d\mu\defeq\int_Eu\,d\mu+i\int_Ev\,d\mu$.
\item \hlt{$L^2(\mu)$}: complex measurable $f$ satisfying $\int_X\abs{f}^2\,d\mu<+\infty$, modulo a.e. equality; $\norm{f}_2^2\defeq\int_X\abs{f}^2\,d\mu$.
\item The \hlt{$L^2$ inner product} is $\innerprod[L^2]{f,g}\defeq\int_Xf\overline{g}\,d\mu$; it induces $\norm{f}_2$.
\end{enumerate}
\end{definition}

\begin{proposition} \label{prop:complex-l1-absolute-integrability-and-dct} \hlt{Complex $L^1$ and Dominated Convergence}\\
The following statements hold.
\begin{enumerate}[label=(\roman*)]
\setlength{\itemsep}{0pt}
\item For complex measurable $f$ on $E$, $f$ is integrable $\Longleftrightarrow \abs{f}\in\mathscr{L}^1(\mu)$ on $E$; then $\abs{\int_Ef\,d\mu}\leq\int_E\abs{f}\,d\mu$.
\item If $f_n\to f$ a.e. on $E$ and $\abs{f_n}\leq g\in\mathscr{L}^1(\mu)$, then $f$ is integrable and $\int_Ef\,d\mu=\lim\limits_n\int_Ef_n\,d\mu$.
\end{enumerate}
\begin{proof}
\begin{enumerate}[label=(\roman*)]
\setlength{\itemsep}{0pt}
\item[$(i)$] ($\Rightarrow$) Write $f=u+iv$. If $u,v\in\mathscr{L}^1$, then $\abs f\leq\abs u+\abs v$, so $\abs f\in\mathscr{L}^1$.
\item[$(i)$] ($\Leftarrow$) If $\abs f\in\mathscr{L}^1$, then $\abs u,\abs v\leq\abs f$, so $u,v\in\mathscr{L}^1$ and $f$ is integrable.\\
If $\int_Ef\,d\mu=re^{i\theta}$ and $h=e^{-i\theta}f$, then $\abs{\int_Ef\,d\mu}=r=\operatorname{Re}\int_Eh\,d\mu=\int_E\operatorname{Re}h\,d\mu\leq\int_E\abs h\,d\mu=\int_E\abs f\,d\mu$.
\item[$(ii)$] Write $f_n=u_n+iv_n$ and $f=u+iv$. Since $\abs{u_n},\abs{v_n}\leq g$ and $u_n\to u$, $v_n\to v$ a.e., real DCT gives $\int_Eu_n\,d\mu\to\int_Eu\,d\mu$ and $\int_Ev_n\,d\mu\to\int_Ev\,d\mu$; hence the complex identity.
\end{enumerate}
\hlt{Structural Commentary}: Complex integration is two real integrals.\\
\hlt{Key Insight}: Rotate the integral to make its value real and nonnegative before applying the real inequality.
\end{proof}
\end{proposition}

\begin{theorem} \label{thm:l2-schwarz-minkowski-and-completeness} \hlt{$L^2$ Geometry}\\
For $f,g\in L^2(\mu)$, the following hold.
\begin{enumerate}[label=(\roman*)]
\setlength{\itemsep}{0pt}
\item \hlt{Schwarz}: $fg\in L^1(\mu)$ and $\int_X\abs{fg}\,d\mu\leq\norm{f}_2\norm{g}_2$.
\item \hlt{Minkowski}: $f+g\in L^2(\mu)$ and $\norm{f+g}_2\leq\norm{f}_2+\norm{g}_2$.
\item \hlt{Completeness}: modulo equality a.e., $(L^2(\mu),\norm{\cdot}_2)$ is complete.
\end{enumerate}
\begin{proof}
\begin{enumerate}[label=(\roman*)]
\setlength{\itemsep}{0pt}
\item Since $2\abs{fg}\leq\abs f^2+\abs g^2$, $fg\in L^1(\mu)$.\\
For $\lambda\in\mathbb R$, $0\leq Q(\lambda)\defeq\int_X(\abs f+\lambda\abs g)^2\,d\mu=\norm f_2^2+2\lambda\int_X\abs{fg}\,d\mu+\lambda^2\norm g_2^2$; hence $\Delta(Q)\leq0$.
\item By (i), $\norm{f+g}_2^2\leq\int_X(\abs f+\abs g)^2\,d\mu\leq(\norm f_2+\norm g_2)^2$.
\item Let $(f_n)$ be Cauchy. Choose $n_k\uparrow\infty$ with $\norm{f_{n_{k+1}}-f_{n_k}}_2\leq2^{-k}$, and set $h_N\defeq\sum\limits_{k<N}\abs{f_{n_{k+1}}-f_{n_k}}$.\\
Then $\norm{h_N}_2\leq\sum\limits_{k<N}2^{-k}$, so Fatou gives $h\defeq\lim\limits_Nh_N\in L^2$ and $f_{n_k}\to f$ a.e. for some $f$.\\
For fixed $n$, Fatou gives $\norm{f-f_n}_2\leq\liminf\limits_k\norm{f_{n_k}-f_n}_2$; the Cauchy property sends the right side to $0$.
\end{enumerate}
\hlt{Structural Commentary}: Schwarz and Minkowski supply geometry; rapid subsequences prove completeness.\\
\hlt{Key Insight}: Positivity of the integral transfers Euclidean quadratic estimates to function classes.
\end{proof}
\end{theorem}

\begin{example} \label{ex:mazur-l1-l2-homeomorphism} \hlt{Mazur Homeomorphism Between $L^1$ and $L^2$}\\
Let $E\subseteq\R$ be Lebesgue measurable. Define $\Phi(f)\defeq\operatorname{sgn}(f)\abs f^{\frac{1}{2}}$ on $L^1(E)$ and $\Psi(g)\defeq\operatorname{sgn}(g)\abs g^2$ on $L^2(E)$.\\
Then $\Phi:L^1(E)\to L^2(E)$ is a homeomorphism with inverse $\Psi:L^2(E)\to L^1(E)$.
\begin{proof}
$\abs{\Phi(f)}^2=\abs f$ and $\abs{\Psi(g)}=\abs g^2$, so the maps are well-defined; pointwise $\Psi\Phi=\mathrm{id}_{L^1}$ and $\Phi\Psi=\mathrm{id}_{L^2}$.\\
For $a,b\in\R$, $\abs{\operatorname{sgn}(a)\abs a^{\frac{1}{2}}-\operatorname{sgn}(b)\abs b^{\frac{1}{2}}}^2\leq2\abs{a-b}$, hence $\norm{\Phi(f)-\Phi(h)}_2^2\leq2\norm{f-h}_1$.\\
Also $\abs{\operatorname{sgn}(a)\abs a^2-\operatorname{sgn}(b)\abs b^2}\leq2\abs{a-b}(\abs a+\abs b)$, so $\norm{\Psi(g)-\Psi(k)}_1\leq2\norm{g-k}_2(\norm g_2+\norm k_2)$ by Cauchy--Schwarz. Thus $\Phi$ and $\Psi$ are continuous.\\
\hlt{Key Insight}: Nonlinear power maps can homeomorphically identify different $L^p$ norm topologies.
\end{proof}
\end{example}

\begin{lemma} \label{lem:young-inequality-for-holder} \hlt{Young's Inequality}\\
If $A,B\geq0$, $1<p<\infty$, and $p^{-1}+q^{-1}=1$, then:
\begin{enumerate}[label=(\roman*)]
\setlength{\itemsep}{0pt}
\item $AB\leq\frac{A^p}{p}+\frac{B^q}{q}$;
\item equality holds iff $A^p=B^q$.
\end{enumerate}
\begin{proof}
\begin{enumerate}[label=(\roman*)]
\setlength{\itemsep}{0pt}
\item If $A,B>0$, set $\alpha\defeq\frac{1}{p}$ and $t\defeq\frac{A^p}{B^q}$. Since $q^{-1}=1-\alpha$, weighted AM--GM gives $t^\alpha\leq\alpha t+1-\alpha$.\\
Multiplying by $B^q$ yields $AB\leq\frac{A^p}{p}+\frac{B^q}{q}$. If $AB=0$, the inequality is immediate.
\item For $A,B>0$, equality in weighted AM--GM is $t=1$, i.e. $A^p=B^q$.\\
If $AB=0$, equality in (i) holds iff $A=B=0$, again equivalent to $A^p=B^q$.
\end{enumerate}
\hlt{Structural Commentary}: Young's inequality is weighted AM--GM in conjugate exponents.\\
\hlt{Key Insight}: Rescale to $A^p$ and $B^q$, then use convexity.
\end{proof}
\end{lemma}

\begin{theorem} \label{thm:holder-inequality-lp} \hlt{H\"older's Inequality}\\
Let $1<p<\infty$, $p^{-1}+q^{-1}=1$, $f\in L^p(\mu)$, and $g\in L^q(\mu)$.
\begin{enumerate}[label=(\roman*)]
\setlength{\itemsep}{0pt}
\item $fg\in L^1(\mu)$.
\item $\norm{fg}_1\leq\norm f_p\norm g_q$.
\end{enumerate}
\begin{proof} If $\norm f_p\norm g_q=0$, then $f=0$ a.e. or $g=0$ a.e., so both claims hold.\\
Assume $\norm f_p\norm g_q>0$ and set $A\defeq\frac{\abs f}{\norm f_p}$, $B\defeq\frac{\abs g}{\norm g_q}$. Young gives $\frac{\abs{fg}}{\norm f_p\norm g_q}\leq\frac{A^p}{p}+\frac{B^q}{q}$.
\begin{enumerate}[label=(\roman*)]
\setlength{\itemsep}{0pt}
\item Since $\int A^p\,d\mu=\int B^q\,d\mu=1$, the majorant $\frac{A^p}{p}+\frac{B^q}{q}$ lies in $L^1(\mu)$; hence $fg\in L^1(\mu)$.
\item Integrating the displayed inequality gives $\frac{\norm{fg}_1}{\norm f_p\norm g_q}\leq\frac{1}{p}\int A^p\,d\mu+\frac{1}{q}\int B^q\,d\mu=p^{-1}+q^{-1}=1$.
\end{enumerate}
\hlt{Structural Commentary}: Normalize both factors to unit $L^p$ and $L^q$ size.\\
\hlt{Key Insight}: Products are controlled by conjugate powers.
\end{proof}
\end{theorem}

\begin{definition} \label{def:conjugate-indices-lp} \hlt{Conjugate Indices}\\
For $p,q\in(1,\infty)$, $p$ and $q$ are \hlt{conjugate indices} if $p^{-1}+q^{-1}=1$, equivalently $p+q=pq$.\\
The unique self-conjugate index is $p=q=2$.
\end{definition}

\begin{corollary} \label{cor:cauchy-schwarz-l2-integral} \hlt{Cauchy--Bunyakovsky--Schwarz in $L^2$}\\
If $f,g\in L^2(\mu)$, then:
\begin{enumerate}[label=(\roman*)]
\setlength{\itemsep}{0pt}
\item $fg\in L^1(\mu)$;
\item $\abs{\int_Xfg\,d\mu}\leq\int_X\abs{fg}\,d\mu\leq\norm f_2\norm g_2$.
\end{enumerate}
\begin{proof}
\begin{enumerate}[label=(\roman*)]
\setlength{\itemsep}{0pt}
\item Since $2^{-1}+2^{-1}=1$, Theorem \ref{thm:holder-inequality-lp}(i) gives $fg\in L^1(\mu)$.
\item Theorem \ref{thm:holder-inequality-lp}(ii) gives $\int_X\abs{fg}\,d\mu=\norm{fg}_1\leq\norm f_2\norm g_2$, and $\abs{\int_Xfg\,d\mu}\leq\int_X\abs{fg}\,d\mu$.
\end{enumerate}
\hlt{Structural Commentary}: H\"older at $p=q=2$ is Cauchy--Schwarz.\\
\hlt{Key Insight}: Pair two $L^2$ factors through $2^{-1}+2^{-1}=1$.
\end{proof}
\end{corollary}

\begin{theorem} \label{thm:minkowski-inequality-lp} \hlt{Minkowski's Inequality}\\
If $1\leq p<\infty$ and $f,h\in L^p(\mu)$, then:
\begin{enumerate}[label=(\roman*)]
\setlength{\itemsep}{0pt}
\item $f+h\in L^p(\mu)$;
\item $\norm{f+h}_p\leq\norm f_p+\norm h_p$.
\end{enumerate}
\begin{proof}
For $p=1$, integrate $\abs{f+h}\leq\abs f+\abs h$. Let $p>1$ and $q\defeq\frac{p}{p-1}$.
\begin{enumerate}[label=(\roman*)]
\setlength{\itemsep}{0pt}
\item Since $\abs{f+h}^p\leq2^{p-1}(\abs f^p+\abs h^p)$, $f+h\in L^p(\mu)$. Also $\abs{f+h}^{p-1}\in L^q(\mu)$ because $(p-1)q=p$.
\item Pointwise, $\abs{f+h}^p\leq(\abs f+\abs h)\abs{f+h}^{p-1}$.\\
By H\"older, $\norm{f+h}_p^p\leq(\norm f_p+\norm h_p)\norm{\abs{f+h}^{p-1}}_q=(\norm f_p+\norm h_p)\norm{f+h}_p^{p-1}$.\\
If $\norm{f+h}_p>0$, divide by $\norm{f+h}_p^{p-1}$; otherwise the claim is trivial.
\end{enumerate}
\hlt{Structural Commentary}: H\"older turns the pointwise triangle inequality into a norm triangle inequality.\\
\hlt{Key Insight}: Minkowski is exactly the closure and triangle estimate for $L^p$.
\end{proof}
\end{theorem}

\begin{theorem} \label{thm:lp-normed-linear-space} \hlt{$L^p$ Is Normed}\\
For $1\leq p<\infty$, $(L^p(\mu),\norm{\cdot}_p)$ is a normed linear space.
\begin{proof}
\begin{enumerate}[label=(\roman*)]
\setlength{\itemsep}{0pt}
\item If $f=f'$ and $g=g'$ a.e., then $\alpha f+\beta g=\alpha f'+\beta g'$ a.e..\\
Hence pointwise operations descend to $L^p(\mu)$. Minkowski gives $\alpha f+\beta g\in L^p(\mu)$.
\item $\norm{\alpha f}_p^p=\int_X\abs{\alpha f}^p\,d\mu=\abs{\alpha}^p\norm f_p^p$, so $\norm{\alpha f}_p=\abs{\alpha}\norm f_p$.
\item $\norm f_p=0\Rightarrow\int_X\abs f^p\,d\mu=0\Rightarrow f=0$ a.e.; hence $[f]=0$ in $L^p(\mu)$.
\item Minkowski gives $\norm{f+g}_p\leq\norm f_p+\norm g_p$.
\end{enumerate}
\hlt{Structural Commentary}: Quotienting by a.e. equality turns pointwise operations into vector-space operations.\\
\hlt{Key Insight}: Minkowski supplies the nontrivial norm axiom.
\end{proof}
\end{theorem}

\begin{definition} \label{def:lp-cauchy-complete-language} \hlt{Completeness Language}\\
Let $(V,\norm{\cdot})$ be a normed space.
\begin{enumerate}[label=(\roman*)]
\setlength{\itemsep}{0pt}
\item \hlt{Complete}: $V$ is complete iff every Cauchy sequence $(x_n)\subset V$ has $x\in V$ with $\norm{x_n-x}\to0$.
\item \hlt{Banach}: A Banach space is a complete normed space.
\item \hlt{$L^p$ criterion}: $L^p(\mu)$ is complete iff every $L^p$-Cauchy sequence converges in $\norm{\cdot}_p$.
\end{enumerate}
\end{definition}

\begin{lemma} \label{lem:lp-convergent-sequences-are-cauchy} \hlt{Convergent Implies Cauchy}\\
If $f_n\to f$ in $L^p(\mu)$, then $(f_n)$ is Cauchy in $L^p(\mu)$.
\begin{proof}
Given $\varepsilon>0$, choose $N$ so $n\geq N$ implies $\norm{f_n-f}_p<\frac{\varepsilon}{2}$.\\
Then $m,n\geq N$ gives $\norm{f_m-f_n}_p\leq\norm{f_m-f}_p+\norm{f-f_n}_p<\varepsilon$ by Minkowski.\\
\hlt{Structural Commentary}: A convergent sequence has uniformly small tails around its limit.\\
\hlt{Key Insight}: Use the common limit as the bridge between two tails.
\end{proof}
\end{lemma}

\begin{definition} \label{def:rapidly-cauchy-sequence} \hlt{Rapidly Cauchy Sequence}\\
Let $(V,\norm{\cdot})$ be a normed space. A sequence $(x_k)\subset V$ is \hlt{rapidly Cauchy} if $\exists(\varepsilon_k)\subset(0,+\infty)$ such that:
\begin{enumerate}[label=(\roman*)]
\setlength{\itemsep}{0pt}
\item $\sum\limits_{k=1}^{\infty}\varepsilon_k<+\infty$.
\item $\norm{x_{k+1}-x_k}\leq\varepsilon_k^2$ for every $k\geq1$.
\end{enumerate}
\end{definition}

\begin{proposition} \label{prop:rapidly-cauchy-subsequence-extraction} \hlt{Rapidly Cauchy Subsequences}\\
In a normed space, the following hold.
\begin{enumerate}[label=(\roman*)]
\setlength{\itemsep}{0pt}
\item Every rapidly Cauchy sequence is Cauchy.
\item Every Cauchy sequence has a rapidly Cauchy subsequence.
\end{enumerate}
\begin{proof}
\begin{enumerate}[label=(\roman*)]
\setlength{\itemsep}{0pt}
\item Let $(x_k)$ be rapidly Cauchy with witnesses $(\varepsilon_k)$. For $m>n$, $\norm{x_m-x_n}\leq\sum\limits_{j=n}^{m-1}\norm{x_{j+1}-x_j}\leq\sum\limits_{j=n}^{\infty}\varepsilon_j^2$.\\
Since $\sum_j\varepsilon_j<+\infty$, eventually $\varepsilon_j\leq1$, so $\sum\limits_{j=n}^{\infty}\varepsilon_j^2\leq\sum\limits_{j=n}^{\infty}\varepsilon_j\to0$.
\item Let $(x_n)$ be Cauchy. Choose $n_1<n_2<\cdots$ such that $m,n\geq n_k\Rightarrow\norm{x_m-x_n}\leq4^{-k}$.\\
Set $y_k\defeq x_{n_k}$ and $\varepsilon_k\defeq2^{-k}$. Then $\sum_k\varepsilon_k<+\infty$ and $\norm{y_{k+1}-y_k}\leq4^{-k}=\varepsilon_k^2$, so $(y_k)$ is rapidly Cauchy.
\end{enumerate}
\hlt{Structural Commentary}: A Cauchy tail contains a summably controlled thread.\\
\hlt{Key Insight}: Replace arbitrary tail control by a geometric subsequence.
\end{proof}
\end{proposition}

\begin{theorem} \label{thm:riesz-fischer-lp-banach} \hlt{Riesz--Fischer Completeness Theorem}\\
For $1\leq p<\infty$, the following hold.
\begin{enumerate}[label=(\roman*)]
\setlength{\itemsep}{0pt}
\item \hlt{Completeness}: $(L^p(\mu),\norm{\cdot}_p)$ is complete.
\item \hlt{A.E. subsequence}: if $f_n\to f$ in $L^p(\mu)$, then some subsequence converges to a representative of $f$ a.e.
\end{enumerate}
\begin{proof}
\begin{enumerate}[label=(\roman*)]
\setlength{\itemsep}{0pt}
\item Let $(f_n)$ be Cauchy. Choose $g_k=f_{n_k}$ with $\norm{g_{k+1}-g_k}_p\leq4^{-k}$. Set $G_N\defeq\abs{g_1}+\sum\limits_{k<N}\abs{g_{k+1}-g_k}$.\\
By Minkowski, $\norm{G_N}_p\leq\norm{g_1}_p+\sum\limits_{k<N}4^{-k}\leq\norm{g_1}_p+\frac{1}{3}$; Fatou gives $G\defeq\lim\limits_NG_N\in L^p$.\\
The series $h\defeq g_1+\sum\limits_k(g_{k+1}-g_k)$ converges a.e. with $\abs h\leq G$ and $g_k\to h$ a.e.\\
Since $\abs{g_k-h}\leq2G$, DCT gives $\norm{g_k-h}_p\to0$. Given $\varepsilon>0$, choose $k$ with $\norm{g_k-h}_p<\frac{\varepsilon}{2}$ and then $N$ with $n\geq N\Rightarrow\norm{f_n-g_k}_p<\frac{\varepsilon}{2}$; hence $\norm{f_n-h}_p\to0$.
\item If $f_n\to f$ in $L^p$, then $(f_n)$ is Cauchy. Applying (i) gives a subsequence $f_{n_k}\to h$ a.e. and in $L^p$.\\
Since also $f_{n_k}\to f$ in $L^p$, uniqueness in normed space gives $[h]=f$ in $L^p$, so $h$ is a representative of $f$.
\end{enumerate}
\hlt{Structural Commentary}: A rapid subsequence converges a.e.; DCT gives norm convergence.\\
\hlt{Key Insight}: Fatou finds the pointwise limit, dominated convergence brings it back to $L^p$.
\end{proof}
\end{theorem}

\begin{proposition} \label{prop:linfty-essential-supremum-characterization} \hlt{Essential Supremum Test}\\
If $f\in L^\infty(\mu)$, then:
\begin{enumerate}[label=(\roman*)]
\setlength{\itemsep}{0pt}
\item $\abs f\leq\norm f_\infty$ a.e.;
\item $A<\norm f_\infty\Rightarrow\mu(\{\abs f>A\})>0$.
\end{enumerate}
\begin{proof}
\begin{enumerate}[label=(\roman*)]
\setlength{\itemsep}{0pt}
\item Let $M\defeq\norm f_\infty$. For each $k$, $N_k\defeq\{\abs f>M+\frac{1}{k}\}$ is null. On $X\setminus\bigcup_kN_k$, $\abs f\leq M+\frac{1}{k}$ $\forall k$, hence $\abs f\leq M$.
\item If $A<M$ and $\mu(\{\abs f>A\})=0$, then $\abs f\leq A$ a.e., so $A$ is an essential upper bound. Contradiction.
\end{enumerate}
\hlt{Structural Commentary}: The essential supremum is the least a.e. upper bound.\\
\hlt{Key Insight}: The sharp bound cannot be lowered on a full-measure set.
\end{proof}
\end{proposition}

\begin{theorem} \label{thm:linfty-banach-space} \hlt{$L^\infty$ Is Banach}\\
$(L^\infty(\mu),\norm{\cdot}_\infty)$ is a Banach space.
\begin{proof}
\begin{enumerate}[label=(\roman*)]
\setlength{\itemsep}{0pt}
\item Pointwise operations descend modulo a.e. equality, and $\norm{\alpha f}_\infty=\abs{\alpha}\norm f_\infty$. Proposition \ref{prop:linfty-essential-supremum-characterization}(i) gives definiteness; outside the union of two null sets, $\abs{f+g}\leq\norm f_\infty+\norm g_\infty$, hence $\norm{f+g}_\infty\leq\norm f_\infty+\norm g_\infty$.
\item Let $(f_n)$ be $L^\infty$-Cauchy. For each $m,n$, choose null $N_{mn}$ with $\abs{f_m-f_n}\leq\norm{f_m-f_n}_\infty$ on $X\setminus N_{mn}$, and set $N\defeq\bigcup\limits_{m,n}N_{mn}$. On $X\setminus N$, $(f_n)$ is uniformly Cauchy, hence converges uniformly to a measurable $f$.
\item The uniform convergence off $N$ gives $\norm{f_n-f}_\infty\to0$.\\
For fixed $n_0$, $\abs f\leq\abs{f_{n_0}}+\sup\limits_{n\geq n_0}\norm{f_n-f_{n_0}}_\infty$ a.e.; hence $f\in L^\infty(\mu)$.
\end{enumerate}
\hlt{Structural Commentary}: Completeness is uniform convergence off one null set.\\
\hlt{Key Insight}: Essential bounds become ordinary uniform bounds off a common null set.
\end{proof}
\end{theorem}


\begin{proposition} \label{prop:simple-density-support-and-tails-lp} \hlt{Simple Approximation, Support, and Tails}\\
Let $1\leq p<\infty$ and $f\in L^p(\mu)$.
\begin{enumerate}[label=(\roman*)]
\setlength{\itemsep}{0pt}
\item Simple measurable functions are dense in $L^p(\mu)$; also in $L^\infty(\mu)$ for $\norm{\cdot}_\infty$.
\item $\{\abs f>0\}$ is $\sigma$-finite.
\item For $E_n\defeq\{\abs f\geq n\}$, $\mu(E_n)\to0$.
\item $\forall\varepsilon>0\ \exists E_\varepsilon\in\mathcal M$ with $\mu(E_\varepsilon)<+\infty$ such that $F\cap E_\varepsilon=\emptyset$ implies $\norm{f\mathbf 1_F}_p<\varepsilon$.
\end{enumerate}
\begin{proof}
\begin{enumerate}[label=(\roman*)]
\setlength{\itemsep}{0pt}
\item Choose simple $s_n\to f$ a.e. with $\abs{s_n}\leq\abs f$ in the real case; treat real and imaginary parts separately in the complex case. Then $\abs{f-s_n}^p\leq2^p\abs f^p$, so DCT gives $\norm{f-s_n}_p\to0$.\\
For $f\in L^\infty$, discard null $N$ with $\abs f\leq\norm f_\infty+\eta$ on $X\setminus N$, then quantize that bounded range.
\item Put $A_k\defeq\{\abs f\geq\frac{1}{k}\}$. Then $\mu(A_k)\leq k^p\norm f_p^p<+\infty$ and $\{\abs f>0\}=\bigcup\limits_kA_k$.
\item $n^p\mu(E_n)\leq\int_{E_n}\abs f^p\,d\mu\leq\norm f_p^p$, hence $\mu(E_n)\leq n^{-p}\norm f_p^p\to0$.
\item With $A_k$ as in (ii), $A_k\uparrow\{\abs f>0\}$ and $\abs f^p\mathbf 1_{X\setminus A_k}\to0$ a.e. DCT gives $\int_{X\setminus A_k}\abs f^p\,d\mu\to0$.\\
Choose $E_\varepsilon=A_k$ with this integral $<\varepsilon^p$; if $F\cap E_\varepsilon=\emptyset$, then $\norm{f\mathbf 1_F}_p^p\leq\int_{X\setminus E_\varepsilon}\abs f^p\,d\mu<\varepsilon^p$.
\end{enumerate}
\hlt{Structural Commentary}: $L^p$ mass is simple-approximable and finite-piece capturable.\\
\hlt{Key Insight}: Integrability turns level sets into finite-measure exhaustions.
\end{proof}
\end{proposition}

\begin{proposition} \label{prop:lp-layer-shell-criterion-finite-measure} \hlt{Shell Criterion on Finite Measure Spaces}\\
Let $\mu(X)<+\infty$, $f$ be measurable, and $E_n\defeq\{n-1\leq\abs f<n\}$ for $n\geq1$.\\
For $1\leq p<\infty$, $f\in L^p(\mu)\iff\sum\limits_{n\geq1}n^p\mu(E_n)<+\infty$.
\begin{proof}
\begin{enumerate}[label=\roman*.]
\setlength{\itemsep}{0pt}
\item[$(\Rightarrow)$] If $f\in L^p(\mu)$, then $\mu(E_1)\leq\mu(X)<+\infty$ and, for $n\geq2$, $n^p\mathbf{1}_{E_n}\leq2^p\abs f^p\mathbf{1}_{E_n}$.\\
Thus $\sum\limits_n n^p\mu(E_n)\leq\mu(X)+2^p\int_X\abs f^p\,d\mu<+\infty$.
\item[$(\Leftarrow)$] If $\sum\limits_n n^p\mu(E_n)<+\infty$, then $\abs f^p=\sum\limits_n\abs f^p\mathbf{1}_{E_n}\leq\sum\limits_n n^p\mathbf{1}_{E_n}$, hence $\int_X\abs f^p\,d\mu<+\infty$.
\end{enumerate}
\hlt{Structural Commentary}: The finite base shell handles small values; weighted shells control tails.\\
\hlt{Key Insight}: Shell weights characterize $L^p$ mass.
\end{proof}
\end{proposition}

\begin{proposition} \label{prop:lp-inclusions-interpolation-and-multipliers} \hlt{Inclusions, Interpolation, and Multipliers}
\begin{enumerate}[label=(\roman*)]
\setlength{\itemsep}{0pt}
\item If $\mu(X)<+\infty$ and $1\leq r\leq p<\infty$, then $L^p(\mu)\subseteq L^r(\mu)$ and $\norm f_r\leq\norm f_p\mu(X)^{\frac{1}{r}-\frac{1}{p}}$.
\item On $(\N,2^\N,\#)$, $1\leq p\leq s<\infty$ implies $\ell^p\subseteq\ell^s$ and $\norm a_s\leq\norm a_p$.
\item If $1\leq p_1<p_2<\infty$ and $f\in L^{p_1}(\mu)\cap L^{p_2}(\mu)$, then $f\in L^p(\mu)$ for $p_1\leq p\leq p_2$.
\item If $1\leq p\leq\infty$, then $L^p(\mu)\cdot L^\infty(\mu)\subseteq L^p(\mu)$ and $\norm{fg}_p\leq\norm f_p\norm g_\infty$.
\item $L^\infty(\mu)\subseteq L^1(\mu)\Longleftrightarrow\mu(X)<+\infty$. If $\mu(X)=1$, then $\norm f_\infty=\lim\limits_{p\to\infty}\norm f_p$ for $f\in L^\infty(\mu)$.
\end{enumerate}
\begin{proof}
\begin{enumerate}[label=(\roman*)]
\setlength{\itemsep}{0pt}
\item For $r=p$, $L^p\subseteq L^p$ and the estimate reads $\norm f_p\leq\norm f_p$, i.e. equality.\\
For $r<p$, H\"older applied to $\abs f^r\cdot1$ gives $\int_X\abs f^r\,d\mu\leq\norm f_p^r\mu(X)^{1-\frac{r}{p}}$; take $r$th roots.
\item For $a\in\ell^p$, $\forall n,\ \abs{a_n}\leq\norm a_p$, hence $\sum\limits_n\abs{a_n}^s\leq\norm a_p^{s-p}\sum\limits_n\abs{a_n}^p\leq\norm a_p^s$.
\item For $p\in[p_1,p_2]$, $\abs f^p\leq\abs f^{p_1}\mathbf 1_{\{\abs f\leq1\}}+\abs f^{p_2}\mathbf 1_{\{\abs f>1\}}$.
\item For $p<\infty$, $\abs{fg}^p\leq\norm g_\infty^p\abs f^p$ a.e.; for $p=\infty$, $\abs{fg}\leq\norm f_\infty\norm g_\infty$ a.e.
\item[$(v)$]($\Rightarrow$) If $L^\infty(\mu)\subseteq L^1(\mu)$, then $\mathbf 1_X\in L^\infty(\mu)$ lies in $L^1(\mu)$, so $\mu(X)=\int_X\mathbf 1_X\,d\mu<+\infty$.
\item[$(v)$]($\Leftarrow$) If $\mu(X)<+\infty$ and $f\in L^\infty(\mu)$, then $\int_X\abs f\,d\mu\leq\norm f_\infty\mu(X)<+\infty$. If $\mu(X)=1$, then $\norm f_p\leq\norm f_\infty$.\\
For $A<\norm f_\infty$, $\mu(\{\abs f>A\})>0$ and $\norm f_p\geq A\mu(\{\abs f>A\})^{\frac{1}{p}}\to A$. Let $A\uparrow\norm f_\infty$.
\end{enumerate}
\hlt{Structural Commentary}: Inclusions depend on domain size and pointwise bounds.\\
\hlt{Key Insight}: Check whether size comes from the function values or from the measure of the domain.
\end{proof}
\end{proposition}

\begin{example} \label{ex:lp-membership-models-from-exercises} \hlt{Membership Models}
\begin{enumerate}[label=(\roman*)]
\setlength{\itemsep}{0pt}
\item On $(\N,2^\N,\#)$, $a_n=\frac{1}{n}$ satisfies $\norm a_p^p=\sum\limits_{n\geq1}n^{-p}$; hence $a\in\ell^p\Longleftrightarrow p>1$, while $a\in\ell^\infty\setminus\ell^1$.
\item On $\N$ with $\lambda(\{n\})=n^{-2}$, $\lambda(\N)<+\infty$ and $f(n)=n^{\frac{1}{p_0}}$ satisfies $\int_\N\abs f^p\,d\lambda=\sum\limits_{n\geq1}n^{\frac{p}{p_0}-2}$.\\
Hence $f\in L^p(\lambda)\Longleftrightarrow1\leq p<p_0$. In particular, $f(n)=\sqrt n$ gives $p_0=2$.
\item On $(0,\infty)$, $f(x)=x^{-\frac{1}{2}}(1+\abs{\log x})^{-1}$ gives $\int_0^\infty\abs f^p\,dx=\int_{-\infty}^{\infty}e^{(1-\frac{p}{2})t}(1+\abs t)^{-p}\,dt$ under $x=e^t$.\\
\end{enumerate}
\hlt{Structural Commentary}: $L^p$ membership reduces to series and endpoint tests.\\
\hlt{Key Insight}: Different ends impose different $p$-conditions.
\end{example}

\begin{proposition} \label{prop:holder-dual-norm-attainment} \hlt{H\"older Dual Norm Attainment}\\
Let $1<p<\infty$, $p^{-1}+q^{-1}=1$, and let $f\in L^p(\mu)$ be real-valued.
\begin{enumerate}[label=(\roman*)]
\setlength{\itemsep}{0pt}
\item $\sup\limits_{\norm{g}_q\leq1}\abs{\int_Xfg\,d\mu}=\norm f_p$.
\item If $f\neq0$, the supremum is attained by $g_0=\norm f_p^{-\frac{p}{q}}\sign(f)\abs f^{p-1}$.
\end{enumerate}
\begin{proof}
\begin{enumerate}[label=(\roman*)]
\setlength{\itemsep}{0pt}
\item For $\norm{g}_q\leq1$, H\"older gives $\abs{\int_Xfg\,d\mu}\leq\norm f_p\norm{g}_q\leq\norm f_p$.\\
If $f=0$ this is sharp; if $f\neq0$, item (ii) gives an admissible $g_0$ with $\int_Xfg_0\,d\mu=\norm f_p$.
\item Since $(p-1)q=p$, $\norm{g_0}_q^q=\norm f_p^{-p}\int_X\abs f^p\,d\mu=1$ and $\int_Xfg_0\,d\mu=\norm f_p^{-\frac{p}{q}}\int_X\abs f^p\,d\mu=\norm f_p$.
\end{enumerate}
\hlt{Structural Commentary}: Equality in H\"older is obtained by matching the conjugate powers.\\
\hlt{Key Insight}: The $L^q$ test function points in the power direction of $f$.
\end{proof}
\end{proposition}

\begin{definition} \label{def:bounded-linear-functional-on-lp} \hlt{Bounded Linear Functionals on $L^p$}
\begin{enumerate}[label=(\roman*)]
\setlength{\itemsep}{0pt}
\item A \hlt{linear functional} on real $L^p(\mu)$ is a linear map $G:L^p(\mu)\to\mathbb R$.
\item $G$ is \hlt{bounded} if $\exists M<+\infty$ such that $\forall f\in L^p(\mu)$, $\abs{G(f)}\leq M\norm f_p$.
\item The \hlt{operator norm} is $\norm G\defeq\sup\{\abs{G(f)}:f\in L^p(\mu),\ \norm f_p\leq1\}$.
\item $G$ is \hlt{positive} if $f\geq0$ a.e. implies $G(f)\geq0$.
\end{enumerate}
\end{definition}

\begin{lemma} \label{lem:bounded-lp-functional-positive-split} \hlt{Positive Split of a Bounded Functional}\\
Let $1\leq p<+\infty$ and $G\in(L^p(\mu))^*_\mathbb R$. Then $G=G^+-G^-$ for positive bounded linear functionals $G^\pm$.
\begin{proof}
For $f\geq0$, set $G^+(f)\defeq\sup\{G(g):0\leq g\leq f\}$.\\
This is finite since $\abs{G(g)}\leq\norm G\norm g_p\leq\norm G\norm f_p$, and positive homogeneity is immediate.\\
For $f_1,f_2\geq0$, $G^+(f_1)+G^+(f_2)\leq G^+(f_1+f_2)$; conversely, if $0\leq h\leq f_1+f_2$, set $h_1=(h-f_2)^+$ and $h_2=h-h_1$, so $h=h_1+h_2$ with $0\leq h_j\leq f_j$, hence $G(h)\leq G^+(f_1)+G^+(f_2)$.\\
Thus $G^+$ is additive on the positive cone. Extend by $G^+(f)=G^+(f^+)-G^+(f^-)$; then $G^+$ is positive, linear, and bounded. Let $G^-=G^+-G$. Since $G^+(f)\geq G(f)$ for $f\geq0$, $G^-$ is positive and bounded.\\
\hlt{Structural Commentary}: Function-lattice functionals split into positive parts.\\
\hlt{Key Insight}: Split any competitor $h\leq f_1+f_2$ into the portions assigned to $f_1$ and $f_2$.
\end{proof}
\end{lemma}

\begin{theorem} \label{thm:riesz-representation-l1} \hlt{Riesz Representation for $L^1$}\\
Let $(X,\mathcal M,\mu)$ be $\sigma$-finite and $G\in(L^1(\mu))^*_\mathbb R$.
\begin{enumerate}[label=(\roman*)]
\setlength{\itemsep}{0pt}
\item $\exists g\in L^\infty(\mu)$ such that $\forall f\in L^1(\mu)$, $G(f)=\int_Xfg\,d\mu$.
\item $\norm G=\norm g_\infty$.
\item If $G$ is positive, then $g\geq0$ a.e.
\end{enumerate}
\begin{proof}
\begin{enumerate}[label=(\roman*)]
\setlength{\itemsep}{0pt}
\item First let $\mu(X)<+\infty$ and $G\geq0$. Set $\lambda(E)\defeq G(\mathbf 1_E)$. Finite additivity is linearity; if $E_n\uparrow E$, then $\mathbf 1_{E_n}\to\mathbf 1_E$ in $L^1$, hence $\lambda(E_n)\to\lambda(E)$. Thus $\lambda$ is finite, and $\lambda\ll\mu$ since $\mu(E)=0$ gives $\mathbf 1_E=0$ in $L^1$. Radon--Nikodym gives $g\geq0$ with $\lambda(E)=\int_Eg\,d\mu$, so simple functions are represented.\\
By (ii), $g\in L^\infty$, hence $L^1$-density of simple functions extends the identity to all $f\in L^1$. For $\sigma$-finite $\mu$, apply the finite case to $f\mapsto G(f\mathbf 1_{X_n})$ on $X_n\uparrow X$, $\mu(X_n)<+\infty$, and patch densities on overlaps. For general $G$, use the positive split $G=G^+-G^-$ and subtract the two densities.
\item The pairing estimate gives $\norm G\leq\norm g_\infty$. Conversely, if $A>\norm G$ and $E\subseteq\{\abs g>A\}$ has $\mu(E)<+\infty$, then with $f=\sign(g)\mathbf 1_E$, $A\mu(E)<\int_E\abs g\,d\mu=G(f)\leq\norm G\mu(E)$, so $\mu(E)=0$.\\
A $\sigma$-finite exhaustion gives $\forall A>\norm G,\ \mu(\{\abs g>A\})=0$, hence $\norm g_\infty\leq\norm G$.
\item If $G$ is positive and $B=\{g<0\}$, then $\forall E\subseteq B$ of finite measure, $0\leq G(\mathbf 1_E)=\int_Eg\,d\mu\leq0$ only if $\int_E(-g)\,d\mu=0$. Hence $\mu(E)=0$ on $\{g<-\frac{1}{k}\}$ $\forall k$, and a $\sigma$-finite exhaustion gives $\mu(B)=0$.
\end{enumerate}
\hlt{Structural Commentary}: A bounded functional on $L^1$ is integration against an essentially bounded density.\\
\hlt{Key Insight}: Indicators give a measure; Radon--Nikodym gives the density.
\end{proof}
\end{theorem}

\begin{theorem} \label{thm:riesz-representation-lp} \hlt{Riesz Representation for $L^p$, $1<p<+\infty$}\\
Let $(X,\mathcal M,\mu)$ be any measure space, $1<p<+\infty$, $q=\frac{p}{p-1}$, and $G\in(L^p(\mu))^*_\mathbb R$.
\begin{enumerate}[label=(\roman*)]
\setlength{\itemsep}{0pt}
\item \hlt{Representation}: $\exists g\in L^q(\mu)$ such that $\forall f\in L^p(\mu)$, $G(f)=\int_Xfg\,d\mu$.
\item \hlt{Norm}: every representing $g$ satisfies $\norm G=\norm g_q$.
\item \hlt{Uniqueness}: if $g,h\in L^q(\mu)$ represent $G$, then $g=h$ a.e.
\end{enumerate}
\begin{proof}
\begin{enumerate}[label=(\roman*)]
\setlength{\itemsep}{0pt}
\item Finite $\mu$: for $G\geq0$, set $\lambda(E)\defeq G(\mathbf 1_E)$. Since $E_n\uparrow E\Rightarrow\mathbf 1_{E_n}\to\mathbf 1_E$ in $L^p$, $\lambda$ is finite.\\
Since $\mu(E)=0\Rightarrow\mathbf 1_E=0$ in $L^p$, $\lambda\ll\mu$. Let $g=\frac{d\lambda}{d\mu}$.\\
By simple approximation, $G(\phi)=\int_X\phi g\,d\mu$ for bounded $\phi$.\\
For $g_m\defeq g\wedge m$, $\int_Xg_m^q\,d\mu\leq G(g_m^{q-1})\leq\norm G(\int_Xg_m^q\,d\mu)^{1/p}$; hence $g\in L^q$, by $L^p$-density $G(f)=\int_Xfg\,d\mu$.\\
General finite $G$ follows from $G=G^+-G^-$. For $\sigma$-finite $\mu$, choose $X_n\uparrow X$, $\mu(X_n)<+\infty$, represent $G_n:u\mapsto G(\tilde u)$ by $g_n$, and patch the compatible $g_n$ to $g$.\\
Since $\norm{g_n}_q=\norm{G_n}\leq\norm G$, by MCT $g\in L^q$; $f\mathbf 1_{X_n}\to f$ in $L^p\Rightarrow G(f)=\int_Xfg\,d\mu$.\\
For arbitrary $\mu$, choose $\norm{f_n}_p=1$, $G(f_n)\to\norm G$, and $X_0\defeq\bigcup\limits_{n,k}\{\abs{f_n}>1/k\}$.\\
If $E\subseteq X\setminus X_0$, $\mu(E)<+\infty$, and $a=G(\mathbf 1_E)$, then $G(f_n)\pm ta\leq\norm G(1+t^p\mu(E))^{1/p}$; after $n\to\infty$ and $t\downarrow0$, $a=0$. Thus $G$ vanishes off $X_0$; represent on $X_0$ and extend by $0$.
\item For representing $g$, by H\"older, $\norm G\leq\norm g_q$; by Proposition \ref{prop:holder-dual-norm-attainment}, $\norm g_q=\sup\limits_{\norm f_p\leq1}\abs{\int_Xfg\,d\mu}\leq\norm G$.
\item If $g,h$ represent $G$, then $k\defeq g-h$ satisfies $\int_Xfk\,d\mu=0$ $\forall f\in L^p(\mu)$.\\
By Proposition \ref{prop:holder-dual-norm-attainment} applied to $k\in L^q(\mu)$ with exponents $q,p$, $\norm k_q=0$; hence $g=h$ a.e.
\end{enumerate}
\hlt{Structural Commentary}: For $1<p<\infty$, every bounded functional comes from the conjugate space $L^q$.\\
\hlt{Key Insight}: Localize to $\sigma$-finite support, then patch.
\end{proof}
\end{theorem}

\begin{definition} \label{def:bfa-measures-linfty-dual} \hlt{Bounded Finitely Additive Measures}\\
Let $(X,\mathcal M,\mu)$ be a measure space.\\
For $E\in\mathcal M$, let $\Pi(E)$ be the finite families $(E_1,\ldots,E_n)$ with $E_k\in\mathcal M$, $E_k\subseteq E$, and $E_i\cap E_j=\emptyset$ for $i\neq j$.\\
For $\nu:\mathcal M\to\mathbb R$, set $\abs{\nu}_{\mathrm{fa}}(E)\defeq\sup_{(E_k)\in\Pi(E)}\sum\limits_k\abs{\nu(E_k)}$ and $\norm{\nu}_{\mathrm{var}}\defeq\abs{\nu}_{\mathrm{fa}}(X)$.\\
Define $\operatorname{BFA}(X,\mathcal M,\mu)$ as the set of $\nu:\mathcal M\to\mathbb R$ satisfying:
\begin{enumerate}[label=(\roman*)]
\setlength{\itemsep}{0pt}
\item \hlt{Finite additivity}: $\nu(\emptyset)=0$ and $\nu(\bigsqcup\limits_{k=1}^nE_k)=\sum\limits_{k=1}^n\nu(E_k)$.
\item \hlt{Null compatibility}: $\mu(E)=0\Rightarrow\nu(E)=0$.
\item \hlt{Bounded variation}: $\norm{\nu}_{\mathrm{var}}<+\infty$.
\end{enumerate}
\hlt{Key Insight}: $\operatorname{BFA}$ requires finite additivity, null compatibility, and finite variation, not countable additivity.
\end{definition}

\begin{theorem} \label{thm:kantorovitch-representation-linfty-dual} \hlt{Kantorovitch Representation for $(L^\infty)^*$}\\
Let $(X,\mathcal M,\mu)$ be a measure space.\\
For $\nu\in\operatorname{BFA}(X,\mathcal M,\mu)$ and simple $\varphi=\sum\limits_{k=1}^nc_k\mathbf 1_{E_k}$ over a finite measurable partition, set $I_\nu(\varphi)\defeq\sum\limits_{k=1}^nc_k\nu(E_k)$.
\begin{enumerate}[label=(\roman*)]
\setlength{\itemsep}{0pt}
\item \hlt{Extension}: $I_\nu$ descends uniquely to $T_\nu\in(L^\infty(\mu))^*_\mathbb R$.
\item \hlt{Isometry}: $\norm{T_\nu}=\norm{\nu}_{\mathrm{var}}$.
\item \hlt{Surjectivity}: if $S\in(L^\infty(\mu))^*_\mathbb R$ and $\nu_S(E)\defeq S(\mathbf 1_E)$, then $\nu_S\in\operatorname{BFA}(X,\mathcal M,\mu)$ and $S=T_{\nu_S}$.
\end{enumerate}
Thus $\nu\mapsto T_\nu$ is an isometric isomorphism $\operatorname{BFA}(X,\mathcal M,\mu)\cong(L^\infty(\mu))^*_\mathbb R$.
\begin{proof}
\begin{enumerate}[label=(\roman*)]
\setlength{\itemsep}{0pt}
\item Refinement and finite additivity make $I_\nu$ well defined.\\
For simple $\varphi$, discard $N=\{\abs{\varphi}>\norm{\varphi}_\infty\}$; since $\abs{\nu}_{\mathrm{fa}}(N)=0$, $\abs{I_\nu(\varphi)}\leq\norm{\varphi}_\infty\norm{\nu}_{\mathrm{var}}$.\\
Hence uniform simple approximation extends $I_\nu$ to bounded measurable $f$.\\
If $\operatorname{supp}f\subseteq N$ and $\mu(N)=0$, then $\abs{\nu}_{\mathrm{fa}}(N)=0$, so $I_\nu(f)=0$; therefore the extension depends only on the $L^\infty$ class and defines $T_\nu$ with $\norm{T_\nu}\leq\norm{\nu}_{\mathrm{var}}$.
\item For $(E_k)\in\Pi(X)$, set $c_k\defeq\sign(\nu(E_k))$ and $\varphi\defeq\sum\limits_kc_k\mathbf 1_{E_k}$.\\
Then $\norm{\varphi}_\infty\leq1$ and $\sum\limits_k\abs{\nu(E_k)}=T_\nu(\varphi)\leq\norm{T_\nu}$.\\
Supremizing over $\Pi(X)$ gives $\norm{\nu}_{\mathrm{var}}\leq\norm{T_\nu}$, and item (i) gives equality.
\item For $S\in(L^\infty(\mu))^*_\mathbb R$, $\nu_S(\emptyset)=0$ and finite additivity follow from linearity.\\
If $\mu(E)=0$, then $\mathbf 1_E=0$ in $L^\infty$, so $\nu_S(E)=0$.\\
For $(E_k)\in\Pi(X)$, with $c_k=\sign(\nu_S(E_k))$ and $\varphi=\sum\limits_kc_k\mathbf 1_{E_k}$, $\sum\limits_k\abs{\nu_S(E_k)}=S(\varphi)\leq\norm S$.\\
Hence $\nu_S\in\operatorname{BFA}(X,\mathcal M,\mu)$. Since $S$ and $T_{\nu_S}$ agree on simple functions, uniform simple density gives $S=T_{\nu_S}$.
\end{enumerate}
\hlt{Structural Commentary}: Simple extensions give $T_\nu$; indicators recover $\nu$.\\
\hlt{Key Insight}: Finite additivity supplies the full $L^\infty$ dual.
\end{proof}
\end{theorem}

\begin{proposition} \label{prop:l1-countably-additive-part-of-linfty-dual} \hlt{The Countably Additive Part}\\
For $g\in L^1(\mu)$, set $\nu_g(E)\defeq\int_E g\,d\mu$.
\begin{enumerate}[label=(\roman*)]
\setlength{\itemsep}{0pt}
\item $\nu_g\in\operatorname{BFA}(X,\mathcal M,\mu)$.
\item $\norm{\nu_g}_{\mathrm{var}}=\norm{g}_1$.
\item $T_{\nu_g}(f)=\int_Xfg\,d\mu$ for $f\in L^\infty(\mu)$; hence $L^1(\mu)$ embeds isometrically into $(L^\infty(\mu))^*$.
\end{enumerate}
\begin{proof}
\begin{enumerate}[label=(\roman*)]
\setlength{\itemsep}{0pt}
\item For disjoint $E_1,\ldots,E_n$, $\nu_g(\bigsqcup\limits_kE_k)=\sum\limits_k\nu_g(E_k)$ by finite additivity of the integral, and\\
$\mu(E)=0\Rightarrow\nu_g(E)=0$. Also, for $(E_k)\in\Pi(X)$, $\sum\limits_k\abs{\nu_g(E_k)}\leq\sum\limits_k\int_{E_k}\abs g\,d\mu\leq\norm g_1$, so $\nu_g\in\operatorname{BFA}(X,\mathcal M,\mu)$.
\item Corollary \ref{cor:variation-supremum-formulas-and-null-sets} gives $\norm{\nu_g}_{\mathrm{var}}=\abs{\nu_g}(X)$, and Theorem \ref{thm:integral-charge-variation-density} gives $\abs{\nu_g}(X)=\int_X\abs g\,d\mu=\norm g_1$.
\item If $f=\sum\limits_kc_k\mathbf 1_{E_k}$ is simple, then $T_{\nu_g}(f)=\sum\limits_kc_k\nu_g(E_k)=\int_Xfg\,d\mu$.\\
For $f\in L^\infty(\mu)$, choose simple $f_n\to f$ uniformly; then $T_{\nu_g}(f_n)\to T_{\nu_g}(f)$ and $\int_Xf_ng\,d\mu\to\int_Xfg\,d\mu$.
\end{enumerate}
\hlt{Structural Commentary}: Variation fixes the norm; simple approximation extends the pairing.\\
\hlt{Key Insight}: Countably additive representatives are exactly the familiar $L^1$ integral pairings.
\end{proof}
\end{proposition}

\begin{figure}[H]
  \centering
  \begin{tikzpicture}[
  >=Stealth,
  every node/.style={font=\scriptsize},
  spacebox/.style={draw=blue!55!black, fill=blue!5, rounded corners=1.5pt,
    minimum width=2.10cm, minimum height=0.72cm, align=center,
    inner xsep=6pt, inner ysep=4pt},
  dualbox/.style={draw=violet!55!black, fill=violet!6, rounded corners=1.5pt,
    minimum width=2.35cm, minimum height=0.72cm, align=center,
    inner xsep=6pt, inner ysep=4pt},
  measbox/.style={draw=green!45!black, fill=green!8, rounded corners=1.5pt,
    minimum width=1.80cm, minimum height=0.72cm, align=center,
    inner xsep=6pt, inner ysep=4pt},
  fullbox/.style={draw=red!55!black, fill=red!5, rounded corners=1.5pt,
    minimum width=2.55cm, minimum height=0.72cm, align=center,
    inner xsep=6pt, inner ysep=4pt},
  arrow/.style={->, line width=0.55pt, black!65, shorten <=2pt, shorten >=2pt},
  mapword/.style={font=\tiny, fill=white, inner sep=1.4pt, align=center},
  rowword/.style={font=\tiny\itshape, text=black!68, anchor=east}
]
  \node[rowword] at (-5.70,0.72) {$1<p<\infty$};
  \node[spacebox] (lq) at (-4.25,0.72) {$L^q(\mu)$};
  \node[dualbox] (lpdual) at (4.25,0.72) {$(L^p(\mu))^*$};

  \draw[arrow] (lq) -- node[above, mapword] {$\cong$} (lpdual);

  \node[rowword] at (-5.70,-0.78) {$p=\infty$};
  \node[measbox] (lone) at (-4.25,-0.78) {$L^1(\mu)$};
  \node[fullbox] (bfa) at (0,-0.78) {$\operatorname{BFA}(X,\mathcal M,\mu)$};
  \node[dualbox] (linftydual) at (4.25,-0.78) {$(L^\infty(\mu))^*$};

  \draw[arrow] (lone) -- node[above, mapword] {$g\mapsto\nu_g$} (bfa);
  \draw[arrow] (bfa) -- node[above, mapword] {$\nu\mapsto T_\nu$} (linftydual);
\end{tikzpicture}

  \caption{$L^q(\mu)\cong(L^p(\mu))^*$ for $1<p<\infty$; $L^1(\mu)$ is the countably additive part of $\operatorname{BFA}\cong(L^\infty(\mu))^*$.}
  \label{fig:lp-linfty-duality-map}
\end{figure}

\begin{remark} \label{rem:riesz-l2-route-to-radon-nikodym} \hlt{$L^2$ Riesz Gives Another Radon--Nikodym Route}\\
Let $\lambda,\mu$ be finite and $\lambda\ll\mu$. Put $\rho=\lambda+\mu$ and $G(h)\defeq\int h\,d\lambda$ on $L^2(\rho)$.\\
If $G(h)=\int hg\,d\rho$, indicator tests give $0\leq g\leq1$ $\rho$-a.e. and $\mu(\{g=1\})=0$.\\
Since $\int h(1-g)\,d\lambda=\int hg\,d\mu$, take $E_N=E\cap\{1-g\geq\frac{1}{N}\}$ and $h=\frac{\mathbf 1_{E_N}}{1-g}\in L^2(\rho)$ to get $\lambda(E_N)=\int_{E_N}\frac{g}{1-g}\,d\mu$. Let $N\to\infty$ by monotone convergence and $\lambda\ll\mu$: $\lambda(E)=\int_E \frac{g}{1-g}\,d\mu$.\\
\hlt{Key Insight}: Riesz representation on $L^2(\lambda+\mu)$ encodes the density in the ratio $\frac{g}{1-g}$.
\end{remark}

\begin{lemma} \label{lem:lp-integral-absolute-continuity} \hlt{Absolute Continuity of $L^p$ Mass}\\
For $1\leq p<\infty$ and $f\in L^p(\mu)$, $\forall\varepsilon>0\,\exists\delta>0$ such that $\mu(E)<\delta$ implies $\norm{f\mathbf 1_E}_p<\varepsilon$.
\begin{proof}
Let $h=\abs f^p\in L^1$. Choose $M>0$ with $\int_{\{h>M\}}h\,d\mu<\frac{\varepsilon^p}{2}$, and set $\delta=\frac{\varepsilon^p}{2M}$.\\
If $\mu(E)<\delta$, then $\norm{f\mathbf 1_E}_p^p=\int_Eh\,d\mu\leq M\mu(E)+\int_{\{h>M\}}h\,d\mu<\varepsilon^p$.\\
\hlt{Structural Commentary}: $L^p$ mass is absolutely continuous with respect to measure.\\
\hlt{Key Insight}: Truncate the tail; small measure controls the bounded part.
\end{proof}
\end{lemma}

\begin{proposition} \label{prop:lp-cauchy-local-mass-limits} \hlt{Local Mass Limits for $L^p$ Cauchy Sequences}\\
For $1\leq p<\infty$ and $f_n\in L^p(\mu)$, set $\beta_n(E)\defeq\norm{f_n\mathbf 1_E}_p$ for $E\in\mathcal M$.
The following hold.
\begin{enumerate}[label=(\roman*)]
\setlength{\itemsep}{0pt}
\item $\abs{\beta_n(E)-\beta_m(E)}\leq\norm{f_n-f_m}_p$ for $E\in\mathcal M$.
\item If $(f_n)$ is $L^p$-Cauchy, then $\forall E\in\mathcal M$, $\beta(E)\defeq\lim\limits_n\beta_n(E)$ exists.
\end{enumerate}
\begin{proof}
\begin{enumerate}[label=(\roman*)]
\setlength{\itemsep}{0pt}
\item Since $f_n\mathbf 1_E=f_m\mathbf 1_E+(f_n-f_m)\mathbf 1_E$, $\beta_n(E)-\beta_m(E)\leq\norm{f_n-f_m}_p$.\\
Interchanging $m,n$ gives $\beta_m(E)-\beta_n(E)\leq\norm{f_n-f_m}_p$; hence $\abs{\beta_n(E)-\beta_m(E)}\leq\norm{f_n-f_m}_p$.
\item Fix $E\in\mathcal M$ and $\varepsilon>0$. Choose $N$ such that $m,n\geq N\Rightarrow\norm{f_n-f_m}_p<\varepsilon$.\\
Then item (i) gives $\abs{\beta_n(E)-\beta_m(E)}<\varepsilon$, so $(\beta_n(E))$ converges in $\mathbb R$.
\end{enumerate}
\hlt{Structural Commentary}: Restricting to $E$ is a $1$-Lipschitz operation on $L^p$ mass.\\
\hlt{Key Insight}: Local mass inherits global Cauchy control.
\end{proof}
\end{proposition}

\begin{proposition} \label{prop:lp-cauchy-local-mass-control} \hlt{Uniform Local Mass Control for $L^p$ Cauchy Sequences}\\
Let $1\leq p<\infty$, $f_n$ be $L^p$-Cauchy, and $\beta_n(E)\defeq\norm{f_n\mathbf 1_E}_p$. Then $\forall\varepsilon>0$:
\begin{enumerate}[label=(\roman*)]
\setlength{\itemsep}{0pt}
\item $\exists E_\varepsilon\in\mathcal M$ with $\mu(E_\varepsilon)<+\infty$ such that $F\cap E_\varepsilon=\emptyset$ implies $\sup\limits_n\beta_n(F)<\varepsilon$;
\item $\exists\delta>0$ such that $\mu(E)<\delta$ implies $\sup\limits_n\beta_n(E)<\varepsilon$.
\end{enumerate}
\begin{proof}
Choose $N$ so $n\geq N$ implies $\norm{f_n-f_N}_p<\frac{\varepsilon}{3}$.
\begin{enumerate}[label=(\roman*)]
\setlength{\itemsep}{0pt}
\item Choose finite-measure $E_j$ for $1\leq j\leq N$ with $F\cap E_j=\emptyset\Rightarrow\beta_j(F)<\frac{\varepsilon}{3}$, and set $E_\varepsilon=\bigcup\limits_{j\leq N}E_j$.\\
For $n\leq N$ this works; for $n>N$, $\beta_n(F)\leq\norm{f_n-f_N}_p+\beta_N(F)<\varepsilon$.
\item Choose $\delta_j>0$ for $1\leq j\leq N$ with $\mu(E)<\delta_j\Rightarrow\beta_j(E)<\frac{\varepsilon}{3}$, and set $\delta=\min\limits_{j\leq N}\delta_j$.\\
For $n\leq N$ this works; for $n>N$, $\beta_n(E)\leq\norm{f_n-f_N}_p+\beta_N(E)<\varepsilon$.
\end{enumerate}
\hlt{Structural Commentary}: Cauchy families are uniformly tight and absolutely continuous.\\
\hlt{Key Insight}: Control finitely many early terms directly and the tail by Cauchy closeness.
\end{proof}
\end{proposition}

\begin{method} \label{meth:lebesgue-product-proof-patterns} \hlt{Lebesgue Proof Patterns}
\begin{enumerate}[label=(\roman*)]
\setlength{\itemsep}{0pt}
\item \hlt{Approximate}: for $f\geq0$, use $s\in\mathcal{S}^+(f)$ and $\int_E f\,d\mu=\sup\limits_{s\in\mathcal{S}^+(f)}\mathcal{I}_E(s)$.
\item \hlt{Promote}: prove for $\mathbf{1}_A$, extend to $s=\sum\limits_i c_i\mathbf{1}_{A_i}$, then pass through $s_n\uparrow f$.
\item \hlt{Ignore null sets}: if $f\sim_{\mathrm{a.e.}}g$ on $E$, replace $f$ by $g$ inside integrals over measurable $A\subseteq E$.
\item \hlt{Split signs}: use $f=f^+-f^-$ and $h=u+iv$.
\item \hlt{Exchange limits}: reduce Fatou and dominated convergence to $0\leq f_n\uparrow f$.
\item \hlt{Slice}: prove boxes first; use regularity and null slices to pass to products and Fubini.
\end{enumerate}
\hlt{Key Insight}: Reduce to simple functions, monotone convergence, null sets, and product boxes.
\end{method}
